\section{Contraintes de conception}
\label{secContraintes}

Votre conception devra impérativement respecter les contraintes suivantes.

\begin{UPSTIwarning}{Séparation réseau}
Pour des raisons de \textbf{sécurité}, le robot et l'automate gérant l'IHM
devront se trouver sur un \textbf{réseau séparé} du réseau de l'entreprise.
\end{UPSTIwarning}

\begin{UPSTIinfor}{Adressage existant}
Les ordinateurs de la salle occupent déjà les adresses IP suivantes, sur le
réseau de l'entreprise :
\begin{center}
    \texttt{172.16.180.1/24} à \texttt{172.16.180.17/24}
\end{center}
Ces adresses sont des données imposées : votre plan d'adressage pour le
réseau séparé de la cellule robotique devra en tenir compte et rester
cohérent avec cette contrainte.
\end{UPSTIinfor}

\begin{UPSTIinfor}{Protocole de commande}
La commande de la cellule devra se faire via le protocole
\textbf{EtherNet/IP}. Reportez-vous au document ressource \emph{EtherNet/IP}
fourni en support de cours (ainsi qu'au cours au tableau) pour tous les
éléments théoriques nécessaires : modèle producteur/consommateur, objets CIP,
messagerie implicite/explicite, objets Assembly, etc.
\end{UPSTIinfor}
